\documentclass[UTF8,10pt,a4paper]{ctexart}

\usepackage{fontspec}
\setmainfont{TeX Gyre Termes}
\usepackage{amsmath,amsthm}
\usepackage[final]{microtype}
\usepackage{unicode-math}
\setmathfont{TeX Gyre Termes Math}
\usepackage{amscd}
\usepackage{xcolor}
\usepackage[a4paper,left=30mm,right=30mm,top=27mm,bottom=27mm]{geometry}
\usepackage{hyperref}

\hypersetup{
  unicode=true,
  colorlinks=true,
  linkcolor=blue!45!black,
  citecolor=blue!45!black,
  urlcolor=blue!45!black,
  pdftitle={由 Raynaud 丛构造的 Prill 例外有限平展覆盖},
  pdfsubject={代数曲线的有限平展覆盖与 Raynaud 丛},
  pdfkeywords={Prill 问题, 有限平展覆盖, Raynaud 丛, 有限单值群},
  pdfauthor={FDmd-233}
}

\ctexset{
  section={
    number=\arabic{section},
    format=\large\bfseries\centering,
    beforeskip=1.6ex plus .4ex,
    afterskip=1.2ex plus .2ex
  }
}

\newtheorem{theorem}{定理}[section]
\newtheorem{lemma}[theorem]{引理}
\newtheorem{proposition}[theorem]{命题}
\theoremstyle{remark}
\newtheorem{remark}[theorem]{注}
\renewcommand{\proofname}{证明}

\newcommand{\CC}{\mathbb{C}}
\newcommand{\OO}{\mathcal{O}}
\newcommand{\tens}{\otimes}
\DeclareMathOperator{\Pic}{Pic}
\DeclareMathOperator{\Jac}{Jac}
\DeclareMathOperator{\rk}{rk}
\DeclareMathOperator{\Spec}{Spec}

\title{由 Raynaud 丛构造的 Prill 例外有限平展覆盖}
\author{FDmd-233}
\date{2026 年 7 月 10 日}

\begin{document}

\maketitle

\begin{abstract}
设 $C$ 是亏格 $g\geq2$ 的光滑射影连通复曲线。本文证明，存在连通有限平展
Galois 覆盖 $f\colon X\to C$，使得
\[
  h^0\!\left(X,\OO_X\bigl(f^{-1}(p)\bigr)\right)\geq2
\]
对每个 $p\in C$ 都成立。令 $V=R_g\tens M$ 为秩 $g^g$ 的 Raynaud 丛的一个
行列式归一化扭曲。该向量丛满足 $\deg V=0$、$\det V\cong\OO_C$ 与
$H^0(C,V)=0$，但对所有 $p\in C$ 都有 $H^0(C,V(p))\neq0$。我们证明 $V$
的单值群有限，并选取一个对所有点仍保持上述非消失性质的不可约直和项 $W$。
由 $W$ 的单值表示得到 Galois 覆盖后，对应向量丛 $f_*\OO_X$ 以 $\OO_C$ 与
$W$ 为直和项，从而给出所需的不等式。

\medskip
\noindent\textbf{关键词：}Prill 问题；有限平展覆盖；Raynaud 丛；有限单值群

\noindent\textbf{2020 Mathematics Subject Classification：}
14H30（主）；14H60，14H40（次）
\end{abstract}

\section{引言}

设 $f\colon X\to C$ 是光滑射影连通复曲线之间的有限态射。Prill 提出如下问题：
纤维除子 $f^{-1}(p)$ 能否在一个铅笔线性系中移动？等价地，是否可能对一般点
$p\in C$ 有
\[
  h^0\!\left(X,\OO_X\bigl(f^{-1}(p)\bigr)\right)>1
\]
该问题见 \cite[Chapter~VI, Exercise~D]{ACGH}；对亏格至少为三的曲线，
它也被列为 \cite[Problem~14]{LittProblem14}。Biswas--Butler 与 Mohan Kumar
分别得到了一些部分肯定结果 \cite{BiswasButler,MohanKumar}。Landesman 与
Litt 随后证明：对每条亏格二曲线，都存在次数为 $36$ 的连通有限平展覆盖，
并且上述不等式对每个 $p\in C$ 成立 \cite{LandesmanLitt}。

若一个有限覆盖 $f$ 对所有 $p\in C$ 都满足该不等式，我们称它为
\emph{Prill 例外覆盖}。下面的构造适用于每个 $g\geq2$。

\begin{theorem}\label{thm:main-zh}
设 $C$ 是亏格 $g\geq2$ 的光滑射影连通复曲线。则存在连通有限平展 Galois
覆盖 $f\colon X\to C$，使得
\[
  h^0\!\left(X,\OO_X\bigl(f^{-1}(p)\bigr)\right)\geq2
\]
对每个 $p\in C$ 都成立。
\end{theorem}

当 $g\geq3$ 时，定理~\ref{thm:main-zh} 对
\cite[Problem~14]{LittProblem14} 给出了肯定回答。对有限平坦态射，函数
\[
  p\longmapsto h^0\!\left(X,\OO_X\bigl(f^{-1}(p)\bigr)\right)
\]
上半连续，因此其值至少为二的点构成闭集。如果这个闭集包含 $C$ 的一个稠密
开集，它就等于 $C$。所以“一般纤维”表述与逐点表述等价；下文采用后者，
因为构造直接给出这一形式。

证明使用 Raynaud 构造的两个性质。限制丛 $R_g$ 对每个
$\alpha\in\Pic^0(C)$ 都满足 $H^0(C,R_g\tens\alpha)\neq0$；另一方面，
Jacobian 上的向量丛 $E_g$ 经 $[g]\colon J\to J$ 拉回后成为同一个线丛若干份的直和。我们选取
一个行列式平凡且保持逐点非消失性质的扭曲。在相应的 $C$ 的平展拉回上，重复
出现的线丛是挠线丛，因而扭曲后的向量丛具有有限线性单值群。

\section{Raynaud 丛}

记 $J=\Jac(C)=\Pic^0(C)$。在 $J$ 上固定一个 theta 除子 $\Theta$，并选取
Abel--Jacobi 嵌入 $i\colon C\hookrightarrow J$。下述 Raynaud 构造的形式见
\cite[Section~3]{Raynaud} 与 \cite[Section~2.3]{Beauville}。

\begin{theorem}[Raynaud]\label{thm:raynaud-zh}
对每个整数 $n\geq1$，$J$ 上存在向量丛 $E_n$，满足
\begin{equation}\label{eq:raynaud-pullback-zh}
  [n]^*E_n
  \cong H^0\!\left(J,\OO_J(n\Theta)\right)^*\tens_{\CC}\OO_J(n\Theta).
\end{equation}
其限制 $R_n=i^*E_n$ 半稳定，秩为 $n^g$，斜率为 $g/n$，并且
\begin{equation}\label{eq:raynaud-nonvanishing-zh}
  H^0(C,R_n\tens\alpha)\neq0
  \qquad\text{对每个 }\alpha\in\Pic^0(C).
\end{equation}
\end{theorem}

取 $n=g$，并记
\[
  R:=R_g,
  \qquad
  r:=\rk(R)=g^g.
\]
于是 $R$ 是斜率为一的半稳定向量丛，故 $\deg(R)=r$。

\section{行列式归一化扭曲}

\begin{lemma}\label{lem:twist-zh}
存在 $M\in\Pic^{-1}(C)$，同时满足
\begin{equation}\label{eq:choice-M-zh}
  M^{\tens r}\cong(\det R)^{-1}
  \qquad\text{和}\qquad
  H^0(C,R\tens M)=0.
\end{equation}
\end{lemma}

\begin{proof}
映射
\[
  \Pic^{-1}(C)\longrightarrow\Pic^{-r}(C),
  \qquad
  M\longmapsto M^{\tens r}
\]
是 $J$ 上 $r$ 倍同源映射的一个平移。因此，纤维
\[
  T:=\left\{M\in\Pic^{-1}(C)\mid
  M^{\tens r}\cong(\det R)^{-1}\right\}
\]
恰有 $r^{2g}$ 个元素。

若 $H^0(C,R\tens M)\neq0$，称 $M\in\Pic^{-1}(C)$ 为\emph{坏的}。
任取一个非零截面，可得非零态射 $M^{-1}\to R$。设 $L\subset R$ 是其像的
饱和化，则存在有效除子 $D$ 使 $L\cong M^{-1}(D)$。由于
$\deg(M^{-1})=1$，且 $R$ 是斜率为一的半稳定向量丛，
\[
  1+\deg D=\deg L\leq1.
\]
所以 $D=0$，而 $M^{-1}$ 是 $R$ 的次数一线子丛。

固定一个 Jordan--H\"older 滤过
$0=R_0\subset R_1\subset\cdots\subset R_m=R$。给定次数一线子丛
$L\subset R$，取满足 $L\subset R_k$ 的最小 $k$。诱导态射
$L\to R_k/R_{k-1}$ 非零。其源与目标都是斜率一的稳定向量丛，故该态射为
同构。因此，每个坏的 $M$ 都使 $M^{-1}$ 同构于 $R$ 的某个秩一
Jordan--H\"older 因子。该滤过至多有 $r$ 个因子，所以坏的 $M$ 至多有 $r$
个同构类。

由 $g\geq2$ 与 $r=g^g>1$ 可知 $r^{2g}>r$。因此，$T$ 中至少有一个
$M$ 不是坏的。
\end{proof}

固定引理~\ref{lem:twist-zh} 中的 $M$，并令 $V:=R\tens M$。于是
\begin{equation}\label{eq:V-properties-zh}
  \deg V=0,
  \qquad
  \det V\cong\OO_C,
  \qquad
  H^0(C,V)=0.
\end{equation}
对每个 $p\in C$，线丛 $M(p)$ 的次数为零，故由
\eqref{eq:raynaud-nonvanishing-zh} 得
\begin{equation}\label{eq:V-pointwise-zh}
  H^0(C,V(p))=H^0(C,R\tens M(p))\neq0.
\end{equation}

\section{有限单值群}

\begin{proposition}\label{prop:finite-monodromy-zh}
存在 $C$ 的连通有限平展覆盖，使 $V$ 在其上平凡化。
\end{proposition}

\begin{proof}
作笛卡尔方块
\[
\begin{CD}
  \widetilde C @>{\widetilde i}>> J\\
  @V{q}VV                         @VV{[g]}V\\
  C              @>{i}>>          J.
\end{CD}
\]
态射 $q$ 有限、平展且满射。由 \eqref{eq:raynaud-pullback-zh}，
\[
  q^*V
  \cong
  H^0\!\left(J,\OO_J(g\Theta)\right)^*
  \tens_{\CC}
  \bigl(\widetilde i^*\OO_J(g\Theta)\tens q^*M\bigr).
\]

取 $\widetilde C$ 的一个连通分支 $\widetilde C_0$。它在 $q$ 下的像非空且
既开又闭，所以等于 $C$。记
\[
  N:=\left.
  \bigl(\widetilde i^*\OO_J(g\Theta)\tens q^*M\bigr)
  \right|_{\widetilde C_0}.
\]
因为 $h^0(J,\OO_J(g\Theta))=g^g=r$，$q^*V$ 在 $\widetilde C_0$ 上的限制
为 $N^{\oplus r}$。取行列式，并利用 $\det V\cong\OO_C$，得到
\[
  N^{\tens r}\cong\OO_{\widetilde C_0}.
\]

设 $d$ 是 $N$ 的阶，则 $d\mid r$。选取同构
$N^{\tens d}\cong\OO_{\widetilde C_0}$，
它定义有限平展的 $\mu_d$-torsor（主齐性空间）
\[
  \Spec_{\widetilde C_0}
  \left(\bigoplus_{a=0}^{d-1}N^{-a}\right)
  \longrightarrow\widetilde C_0.
\]
$N$ 在这个 torsor 上的拉回是平凡线丛。其任一连通分支都满射到
$\widetilde C_0$；再与 $\widetilde C_0\to C$ 复合，得到连通有限平展覆盖
$h\colon C'\to C$，且 $h^*V$ 平凡。
\end{proof}

取一个支配 $h$ 的连通有限平展 Galois 覆盖后，可以把 $V$ 看成有限群的一个
复表示所对应的向量丛。

\begin{proposition}\label{prop:irreducible-summand-zh}
存在 $C$ 上非平凡、不可约且具有有限单值群的向量丛 $W$，使得
\[
  H^0(C,W(p))\neq0
  \qquad\text{对每个 }p\in C.
\]
\end{proposition}

\begin{proof}
选取使 $V$ 平凡化的连通有限平展 Galois 覆盖 $u\colon Y\to C$，其 Galois
群记为 $H$。固定平凡化 $u^*V\cong\OO_Y^{\oplus r}$。由于 $Y$ 连通且射影，
$H^0(Y,\OO_Y)=\CC$；所以下降矩阵均为常数，并定义表示
$H\to\operatorname{GL}_r(\CC)$。由 Maschke 定理，
\[
  V\cong\bigoplus_{j=1}^{s}W_j^{\oplus m_j},
\]
各 $W_j$ 都是单值群有限的不可约向量丛，且两两不同构。若其中有平凡丛，
则 $V$ 含有平凡直和项，与 \eqref{eq:V-properties-zh} 矛盾。

对 $1\leq j\leq s$，令
\[
  Z_j:=\{p\in C\mid H^0(C,W_j(p))\neq0\}.
\]
设 $p_1,p_2\colon C\times C\to C$ 为两个投影，$\Delta\subset C\times C$
为对角线。向量丛
$p_1^*W_j\tens\OO_{C\times C}(\Delta)$ 在 $C\times\{p\}$ 上的限制就是
$W_j(p)$。对射影态射 $p_2$ 使用上半连续性，得到每个 $Z_j$ 都是闭集。
由 $V$ 的分解与 \eqref{eq:V-pointwise-zh}，
$C=\bigcup_{j=1}^{s}Z_j$。不可约簇不能写成有限个真闭子集之并，所以某个
$Z_j$ 等于 $C$。
\end{proof}

\section{覆盖的构造}

下文使用的向量丛判别准则是 \cite[Lemma~5.4]{LandesmanLittPW} 的平展情形；
为完整起见，这里写出证明。

\begin{proof}[定理~\ref{thm:main-zh} 的证明]
取命题~\ref{prop:irreducible-summand-zh} 中的 $W$，并设
\[
  \rho\colon\pi_1(C,c)\longrightarrow\operatorname{GL}(W_c)
\]
为其单值表示。像 $G$ 是有限群。子群 $\ker\rho$ 给出连通有限平展 Galois
覆盖 $f\colon X\to C$，其 Galois 群为 $G$。

向量丛 $f_*\OO_X$ 对应于正则 $G$-模 $\CC[G]$。每个 $G$ 的不可约复表示都在
$\CC[G]$ 中出现，重数等于其维数。因此，$\OO_C$ 与 $W$ 是 $f_*\OO_X$ 的
两个不同直和项。

对 $p\in C$，纤维除子是 $p$ 的拉回，故
\[
  \OO_X\bigl(f^{-1}(p)\bigr)\cong f^*\OO_C(p).
\]
由投影公式及 $\OO_C,W\subset f_*\OO_X$，
\begin{align*}
  H^0\!\left(X,\OO_X\bigl(f^{-1}(p)\bigr)\right)
  &\cong H^0\!\left(C,f_*\OO_X\tens\OO_C(p)\right)\\
  &\supset H^0(C,\OO_C(p))\oplus H^0(C,W(p)).
\end{align*}
因为 $p$ 有效，$h^0(C,\OO_C(p))\geq1$；由
命题~\ref{prop:irreducible-summand-zh}，$h^0(C,W(p))\geq1$。所以左端空间对
每个 $p\in C$ 的维数都至少为二。
\end{proof}

\begin{remark}
这个构造不给出有用的覆盖次数上界。起始的 Raynaud 丛秩已为 $g^g$，随后还要
取 Kummer 覆盖与 Galois 闭包。亏格二时，\cite{LandesmanLitt} 中次数 $36$
的覆盖次数小得多。
\end{remark}

\begin{remark}
行列式条件不可缺少。公式 \eqref{eq:raynaud-pullback-zh} 只说明 $R_g$ 的拉回
是同一个线丛若干份的直和。行列式平凡使该线丛成为挠线丛，因而扭曲后的向量
丛的线性单值群有限，而不仅是射影单值群有限。
\end{remark}

\begingroup
\footnotesize
\bibliographystyle{amsplain}
\bibliography{source/references}
\endgroup

\end{document}
